% !TEX root = ../../../tree_notes.tex
% =====================================================================
%  Chapter 14 review -- the 15-station revision contract (chapter_end.txt).
%  Revision-only: tables, one-liners, pointers back via \thmref/\probref/
%  \secref/\coderef/\defref/\exref/\figref.  No derivations.
%  Summit: CF 735E Ostap and Tree -- the supply/demand two-index tree
%  knapsack (state (f,g) at v: nearest black in subtree, farthest uncovered
%  in subtree, both capped by k).  Samples verified live:
%    (n=2,k=0, edge 1-2) -> 1;   (n=2,k=1, edge 1-2) -> 3;
%    (n=4,k=1, path 1-2-3-4) -> 9.
%  Exports prob:c27-summit, lst:c27-summit, lst:c27-kit.
% =====================================================================

\chaptersection{Chapter 14 Summary}

One sentence: \emph{when the DP state stops being a constant tuple and grows
an index $j = 0..V$, the tree DP is still one post-order fold -- but the
child merge becomes a capped convolution, and its cost needs an amortisation
theorem rather than a degree count; the four shapes this chapter drills are
budgeted knapsack-over-children, growing profile states, greedy
leaf-pruning with an exchange proof, and win/lose or Grundy folds on
games.} Chapter~13 folded a fixed tuple through a child loop; Chapter~14
pays for a state that grows, and prices the growth honestly. Every entry
below is a pointer, not a re-derivation.

\begin{center}
\renewcommand{\arraystretch}{1.25}
\begin{tabular}{@{}p{0.24\linewidth} p{0.48\linewidth} p{0.20\linewidth}@{}}
\toprule
\textbf{Result} & \textbf{Statement (with conditions)} & \textbf{Where} \\
\midrule
Budgeted subtree state & $\mathit{dp}[v][j]$ = best value in $v$'s subtree
spending budget $j$; children fold one at a time by a capped convolution
$\mathit{ndp}[i+j] = \max(\mathit{ndp}[i+j],\, \mathit{dp}[v][i] + \mathit{dp}[c][j])$.
& \defref{c27-knap-state} \\
Amortised child-merge & Merging children with inner loops capped by
$\min(\mathit{sz}_{\mathit{done}}, V)$ and $\min(\mathit{sz}_c, V)$ costs
$\bigO{nV}$ \emph{across all nodes}, not $\bigO{n^2}$. Ch.13 previewed; here it
is proved. & \thmref{c27-amortize} \\
Profile DP & When the state is a function of a growing index (length, height,
value cap) rather than a fixed tuple, enumerating the start vertex and running
the profile with a rollback costs $\bigO{n^2 \log n}$ and is provably
sufficient. & \defref{c27-profile}, \thmref{c27-profile-bound} \\
Deepest-uncovered greedy & Root at $1$; sort vertices by decreasing depth;
attach a new edge to the parent of the first still-uncovered vertex. The
coverage of a node's parent is a \emph{superset} of the node's coverage, so
the greedy is optimal. & \defref{c27-leaf-prune}, \thmref{c27-exchange} \\
Two-flag win/lose & On a trie, $\mathit{win}[v] = \bigvee_c \neg\mathit{win}[c]$
and $\mathit{lose}[v] = \bigvee_c \neg\mathit{lose}[c]$ decide the $k$-game
match verdict by three cases at the root. & \probref{c27-game-solved} \\
Sprague--Grundy on trees & $g(v) = \operatorname{mex}$ of children's values; disjoint
sums combine by XOR; Green Hackenbush has the closed form $g(v) = \bigoplus_c (g(c)+1)$.
First wins iff root value $\ne 0$. & \defref{c27-grundy}, \thmref{c27-grundy-xor}, \thmref{c27-hackenbush} \\
\bottomrule
\end{tabular}
\end{center}

\medskip
{\small\textbf{Algorithms from this chapter -- one DFS, one growing state, one merge rule.}}
\begin{center}
\small
\renewcommand{\arraystretch}{1.2}
\setlength{\tabcolsep}{3pt}
\begin{tabular}{@{}p{0.19\linewidth} l l p{0.29\linewidth} p{0.155\linewidth}@{}}
\toprule
\textbf{Algorithm} & \textbf{Time} & \textbf{Memory} & \textbf{When to use (deciding phrase)} & \textbf{Code} \\
\midrule
Budgeted child-merge
  & $\bigO{nV}$ & $\bigO{nV}$
  & ``the state has an index $j$ and children must be merged one at a time''
  & \coderef{c27-knap} \\
Three-class knapsack (855C)
  & $\bigO{n^2}$ & $\bigO{n^2}$
  & ``colour assignments counted by budget of type-$k$ vaults''
  & \coderef{c27-855c-soln} \\
Rollback profile DFS
  & $\bigO{n^2\log n}$ & $\bigO{n}$
  & ``best increasing (or LIS-shaped) path from a fixed start''
  & \coderef{c27-profile}, \coderef{c27-490f-soln} \\
Deepest-uncovered leaf-prune
  & $\bigO{n}$ & $\bigO{n}$
  & ``add minimum edges so every leaf within $R$ of $1$''
  & \coderef{c27-greedy}, \coderef{c27-1029e-soln} \\
Two-flag win/lose DP
  & $\bigO{n\Sigma}$ & $\bigO{n\Sigma}$
  & ``trie game, per-move forced, $k$ rounds chained by loser-starts-next''
  & \coderef{c27-455b-soln} \\
Grundy / Hackenbush fold
  & $\bigO{n}$ & $\bigO{n}$
  & ``impartial cutting game on a rooted tree''
  & \coderef{c27-grundy} \\
Supply/demand 2-index DP (Summit)
  & $\bigO{nk^2}$ & $\bigO{nk}$
  & ``choose $\le k$-spread seeds so every node is covered, count colourings''
  & \coderef{c27-summit} \\
\bottomrule
\end{tabular}
\end{center}

\chaptersection{Master Decision Flowchart}

Post-study ladder: walk top to bottom, stop at the first branch that
matches. It decides \emph{which of this chapter's four DP shapes} the
problem is asking for -- and hands off what Chapter~14 does \emph{not}
own: the answer for every root (Chapter~15's rerooting), per-pair
\emph{query} distance (Chapter~16's LCA), and the $\bigO{n\log n}$
small-to-large upgrade (Chapter~17's DSU-on-tree).

\begin{lstlisting}[language={}, basicstyle=\ttfamily\small, numbers=none, frame=single]
  You are handed a rooted tree and asked for a value.  Where does it live?
  |
  |-- state has an INDEX j = 0..V and children must be merged one at a time?
  |     +-- tree knapsack, capped convolution:  [14.1, CF 855C]
  |        cost O(nV) by thm:c27-amortize; NEVER O(deg(v)^2 * V^2).
  |
  |-- state is a growing profile (length / height / LIS / value cap)?
  |     +-- rollback DFS with per-start enumeration:  [14.2, CF 490F]
  |        profile bound thm:c27-profile-bound: O(n^2 log n) suffices.
  |
  |-- you must pick a MINIMAL set (edges, seeds, guards) that COVERS
  |   a tree objective, and an exchange argument is available?
  |     +-- deepest-uncovered leaf-prune greedy:  [14.3, CF 1029E]
  |        prove cov(parent(v)) superset cov(v) (thm:c27-exchange).
  |
  |-- a game with two players on a tree or a trie?
  |     +-- only win/lose? use two flags:  [14.4, CF 455B]
  |        win[v] = OR !win[c], lose[v] = OR !lose[c], leaf = (F,T).
  |     +-- disjoint sub-games composed?  Grundy + XOR:  [14.4]
  |        g(v) = mex of options; branches XOR (thm:c27-grundy-xor).
  |     +-- cutting game on a rooted tree?  Hackenbush:
  |        g(v) = xor of (g[c] + 1) over children (thm:c27-hackenbush).
  |
  |-- both an OWN supply and a FOREIGN demand in the state?
  |     +-- 2-index supply/demand tree DP:  [Summit, CF 735E]
  |        state (f,g) with f = dist to nearest seed in subtree,
  |        g = dist to farthest uncovered node; cap both by k.
  |
  |-- "for every root"?
  |     +-- NOT this chapter -- Ch.15's rerooting (c20).
  |
  |-- d(u, v) for MANY arbitrary pairs, ONLINE?
  |     +-- NOT this chapter -- Ch.16's LCA arithmetic (c11).
  |
  +-- the O(nV) bound degrades when V = Theta(n)?
        +-- NOT this chapter -- Ch.17's DSU-on-tree (c28) buys O(n log^2 n)
            by re-merging the smaller histogram into the larger.
\end{lstlisting}

\chaptersection{Common Mistakes Checklist}

\begin{itemize}[nosep, leftmargin=1.4em]
  \item \textbf{Inner loop not capped by the running size.} In
        \coderef{c27-knap} the two loops must be bounded by
        $\min(\mathit{sz}_{\mathit{done}}, V)$ and $\min(\mathit{sz}_c, V)$,
        not by $V$ alone. Uncapped loops turn the fold into $\bigO{n^2}$ even
        at $V = \bigO{n}$, because a star root re-scans a fully-zero
        $\mathit{dp}[v][0..V]$ once per child. (\thmref{c27-amortize},
        \probref{c27-knap-solved}).
  \item \textbf{Forgetting $i + j \le V$ in the merge.} The convolution must
        drop pairs whose sum exceeds the budget; writing
        $\mathit{ndp}[i+j] \leftarrow \mathit{ndp}[i+j] \oplus (\cdots)$
        unconditionally is a buffer overflow, and clipping only at the array
        boundary silently accepts illegal spend. (\defref{c27-knap-state}).
  \item \textbf{Merging all children at once.} Ch.13's constant-tuple fold
        can be $\Phi$-combined pairwise; a knapsack state cannot. Merging
        two children simultaneously is a two-index convolution and costs
        $\bigO{V^2}$ per node, defeating \thmref{c27-amortize}. Merge
        children \emph{one at a time}, in any order. (\secref{c27-knapsack}).
  \item \textbf{Running the profile DFS from one root only.} In
        \secref{c27-growing} the profile is not a subtree functional; the
        best path may pass through a node without either endpoint being that
        node. Enumerate \emph{every} start vertex (or every peak) or you lose
        soundness -- \thmref{c27-profile-bound}'s amortisation pays for the
        $\bigO{n}$ starts, not $\bigO{1}$.
  \item \textbf{Greedy exchange that flips the containment direction.} The
        deepest-uncovered greedy in \thmref{c27-exchange} works because
        covering a parent is a \emph{superset} of covering the child. If your
        exchange picks the child instead of the parent, the coverage set
        shrinks and the "optimal" greedy under-covers on stars with $R=2$.
        (\probref{c27-greedy-solved}).
  \item \textbf{Combining branches by $+$ or $\max$ in a Grundy fold.}
        Disjoint sub-games compose by XOR: $x \oplus x = 0$. A node with two
        identical value-$2$ branches is a losing position; $2+2 = 4$ and
        $\max = 2$ both wrongly say "First wins". (\thmref{c27-grundy-xor},
        Trap~1 of \secref{c27-games}).
  \item \textbf{$\operatorname{mex}$ scan too short, or bound by the height.}
        $g(v)$ is at most $\deg(v)$, so $\mathit{seen}$ must have $\deg(v)+1$
        slots; a fixed-size buffer clips legitimate values on high-degree
        nodes. Do not "prove" $g(v) \le \height(v)$ for weighted moves --
        that only holds for pure stalks. (\secref{c27-games}, Trap~2).
  \item \textbf{Omitting the $+1$ in Hackenbush.} Each branch contributes
        $g(c)+1$ -- the $+1$ is the edge you stand on. A leaf has $g=0$
        (empty XOR). Dropping the $+1$ collapses every stalk to zero and
        misreports "Second wins" on trivial single-cut positions.
        (\thmref{c27-hackenbush}).
  \item \textbf{\lstinline{long long} on the knapsack.} CF~855C counts
        assignments modulo $10^9+7$; intermediate products of factorials
        overflow \lstinline{int} at $n \approx 40$. The Summit's DP table
        $\mathit{dp}[v][f][g]$ is indexed mod $10^9+7$ but the loop bounds
        are $\bigO{nk^2}$ and $\sum_{v}$ of that fits $10^8$, so
        \lstinline{long long} is not optional in either. (\probref{c27-knap-solved},
        \probref{c27-summit}).
  \item \textbf{Recursion depth on a legal chain.} Same trap as Ch.13: at
        $n = 10^5$ a chain of recursive \code{knap\_dfs} frames overflows the
        stack. Ch.12's \coderef{c09-iter} two-phase form applies verbatim --
        just carry the $dp[v][\cdot]$ array on the frame instead of a scalar.
\end{itemize}

\chaptersection{WA Clinic}

Verdicts are symptoms; the causes are the short list this chapter drilled.
Diagnose top-down -- every leaf names the training problem.

\begin{figure}[H]
\centering
\begin{tikzpicture}[diagram,
    v/.style={boxnode, fill=primaryblue!12, draw=primaryblue, text width=22mm, align=center},
    b/.style={boxnode, fill=difforange!12, draw=difforange, text width=12mm},
    l/.style={boxnode, fill=successgreen!10, draw=successgreen, text width=58mm,
              font=\scriptsize, align=left}]
  \node[v] (verdict) {\textbf{Verdict on the judge}};
  \node[b, below left=8mm and 6mm of verdict] (wa) {\textbf{WA}};
  \node[b, below right=1mm of verdict] (tle) {\textbf{TLE}};
  \node[b, below left=8mm and 1mm of verdict, xshift=30mm] (re) {\textbf{RE}};
  \node[b, below right=8mm and 6mm of verdict] (mle) {\textbf{MLE}};
  \node[l, below right=1mm and 2mm of wa.south] (wa1)
    {uncapped inner loop mis-counts the budget (misses \thmref{c27-amortize}'s cap; \probref{c27-knap-solved});
     greedy picks the child not the parent (violates \thmref{c27-exchange}'s superset; \probref{c27-greedy-solved});
     Hackenbush omits the $+1$ (\thmref{c27-hackenbush}, \probref{c27-game-solved});
     two-flag verdict ignores $k$-parity (CF~455B's $\mathit{lose}[\text{root}]$ case)};
  \node[l, below right=1mm and 2mm of tle.south] (tle1)
    {tree knapsack with $\bigO{V^2}$ child-merge -- one-at-a-time is required (\defref{c27-knap-state});
     profile DFS from one start instead of all $n$ (\thmref{c27-profile-bound});
     Grundy $\operatorname{mex}$ with an unbounded $\mathit{seen}$ scan};
  \node[l, below right=1mm and 2mm of re.south] (re1)
    {recursive \code{knap\_dfs} on a legal $10^5$ chain -- swap for \coderef{c09-iter}'s two-phase;
     Summit's $f+g$ reset indexing off-by-one at $g=-\infty$ sentinel (\probref{c27-summit})};
  \node[l, below right=1mm and 2mm of mle.south] (mle1)
    {$\mathit{dp}$ as $\bigO{nV}$ matrix with $V \sim n$ unclipped -- cap by $\min(\mathit{sz}_v, V)$;
     Summit's $\bigO{nk^2}$ state at $n=k=100$ still fits, but at $n=3\cdot10^5$ requires Shape~B};
  \draw[flow] (verdict) -- (wa); \draw[flow] (verdict) -- (tle);
  \draw[flow] (verdict) -- (re); \draw[flow] (verdict) -- (mle);
  \draw[flow] (wa) -- (wa1); \draw[flow] (tle) -- (tle1);
  \draw[flow] (re) -- (re1); \draw[flow] (mle) -- (mle1);
\end{tikzpicture}
\caption*{Chapter~14 verdict tree -- only the failures this chapter's problems can produce.}
\end{figure}

\medskip
\noindent Backup table (same content, terse):
\begin{center}
\small
\renewcommand{\arraystretch}{1.2}
\begin{tabular}{@{}llp{0.39\linewidth}l@{}}
\toprule
\textbf{Verdict} & \textbf{Usual cause} & \textbf{One-line fix} & \textbf{Trained in} \\
\midrule
WA  & uncapped inner loop     & \code{lim = min(sz\_done, V); clim = min(sz\_c, V);} & \probref{c27-knap-solved} \\
WA  & greedy takes the child  & \code{attach to parent(v)} (superset coverage)  & \thmref{c27-exchange} \\
WA  & Hackenbush drops $+1$   & add one before XOR-ing the child's $g$-value      & \thmref{c27-hackenbush} \\
WA  & two-flag ignores $k$    & \code{else verdict = (k \% 2) ? First : Second;}   & \probref{c27-game-solved} \\
TLE & $\bigO{V^2}$ merge      & one-child-at-a-time convolution                    & \coderef{c27-knap} \\
TLE & one-start profile       & loop over every start vertex                       & \thmref{c27-profile-bound} \\
RE  & recursion on a chain    & \coderef{c09-iter} two-phase stack                  & \secref{c27-knapsack} \\
\bottomrule
\end{tabular}
\end{center}

\chaptersection{Quick Reference Card}

\begin{quickref}[Chapter 14 at a glance]
\textbf{The state grew; the DFS did not.} Ch.13's \coderef{c09-dfs3}
skeleton returns verbatim -- what changes is that $\mathcal S(v)$ is an
array (or a pair, or a Boolean) indexed by a numeric parameter, and the
combiner $\Phi$ becomes a \emph{merge loop} rather than a per-coordinate
function (\defref{c27-knap-state}, \thmref{c27-amortize}).

\medskip
\textbf{Amortisation is the owed theorem.} Merging children one at a time
with inner loops capped by the running sizes costs $\bigO{nV}$ total
(\thmref{c27-amortize}), proved by charging each iteration to a pair
$(x,y)$ with $\mathrm{LCA}(x,y) = v$ and each pair meeting exactly once.
The cap is what makes the bound linear in $n$ and not quadratic.

\medskip
\textbf{Greedy correctness is a superset argument.} Deepest-uncovered
leaf-prune wins because $\mathrm{cov}(\mathrm{parent}(v)) \supseteq
\mathrm{cov}(v)$ -- swapping "add edge at $v$" to "add edge at
$\mathrm{parent}(v)$" never uncovers anything already covered
(\thmref{c27-exchange}). If the containment flips, the greedy fails.

\medskip
\textbf{Games: two flags cheaper than a Grundy number.} When you only
need win/lose, drop to $\mathit{win}[v] = \bigvee_c \neg\mathit{win}[c]$
and $\mathit{lose}[v] = \bigvee_c \neg\mathit{lose}[c]$. Reach for
$\operatorname{mex}$ and XOR only when positions must be \emph{added}
as disjoint sums (\thmref{c27-grundy-xor}, \thmref{c27-hackenbush}).

\medskip
\textbf{Know your boundaries.} Answer for every root is Ch.15 (\chref{c20});
per-pair query distance is Ch.16 (\chref{c11}); $\bigO{n\log n}$
small-to-large is Ch.17 (\chref{c28}); centroid decomposition is Ch.24
(\chref{c18}). This chapter owns \emph{budgeted and game tree DP with
growing states}.
\end{quickref}

\begin{lstlisting}[caption={Chapter 14 code kit: the capped child-merge, the rollback profile DFS, the deepest-uncovered greedy, the Grundy / Hackenbush fold, the two-flag trie DP, and the Summit's supply/demand 2-index DP}, label={lst:c27-kit}]
#include <bits/stdc++.h>
using namespace std;
int n, V;  vector<vector<int>> adj;

// (1) Budgeted child-merge (thm:c27-amortize):  dp[v][j], children one at a time.   [14.1]  lst:c27-knap
vector<int> sz;  vector<vector<long long>> dp;
void knap(int v, int p) {
    sz[v] = 1;
    dp[v].assign(V + 1, LLONG_MIN / 4);
    dp[v][0] = 0;                                     // seed
    for (int c : adj[v]) if (c != p) {
        knap(c, v);
        int lim  = min(sz[v], V);                     // capped loops -- the theorem
        int clim = min(sz[c], V);
        vector<long long> ndp(V + 1, LLONG_MIN / 4);
        for (int i = 0; i <= lim; ++i) if (dp[v][i] > LLONG_MIN / 4)
            for (int j = 0; j <= clim && i + j <= V; ++j)
                ndp[i + j] = max(ndp[i + j], dp[v][i] + dp[c][j]);
        dp[v].swap(ndp);
        sz[v] += sz[c];
    }
}

// (2) Rollback profile DFS:  extend on entry, UNDO the exact change on exit.         [14.2]  lst:c27-profile
vector<int> tail;                                     // patience tails (monotone array)
int profile_dfs(int v, int p, int w) {                // w = value at v
    int pos = lower_bound(tail.begin(), tail.end(), w) - tail.begin();
    int back = (pos < (int)tail.size()) ? tail[pos] : -1;
    if (pos == (int)tail.size()) tail.push_back(w); else tail[pos] = w;
    int best = (int)tail.size();
    for (int u : adj[v]) if (u != p) best = max(best, profile_dfs(u, v, /* ... */));
    if (back == -1) tail.pop_back(); else tail[pos] = back;   // exact UNDO
    return best;
}

// (3) Deepest-uncovered leaf-prune greedy (thm:c27-exchange).                        [14.3]  lst:c27-greedy
//   order vertices by decreasing depth; on the first still-uncovered v,
//   attach to parent(v), then mark everything within R of parent(v) covered.
//   Answer is the number of attachments.

// (4a) Grundy by mex over options (generic impartial game).                          [14.4]  lst:c27-grundy
int grundy_mex(int v, int p) {
    vector<char> seen(adj[v].size() + 1, 0);
    for (int u : adj[v]) if (u != p) seen[grundy_mex(u, v)] = 1;
    int g = 0; while (seen[g]) ++g; return g;
}
// (4b) Hackenbush closed form: g(v) = XOR over children of (g[c] + 1).
int hackenbush(int v, int p) {
    int g = 0;
    for (int u : adj[v]) if (u != p) g ^= (hackenbush(u, v) + 1);
    return g;
}

// (5) Two-flag win/lose on a trie (CF 455B):  leaf = (win=false, lose=true).        [14.4 solved]
void fold_flags(int v) {
    if (trie[v].empty()) { win[v] = false; lose[v] = true; return; }
    win[v] = false; lose[v] = false;
    for (int u : trie[v]) {
        fold_flags(u);
        win[v]  |= !win[u];
        lose[v] |= !lose[u];
    }
}
// verdict:  !win[root] -> Second;  lose[root] -> First;  else First iff k odd.

// (6) Summit supply/demand 2-index DP (thm:c27-amortize promoted to (f,g)-budget).  [Summit 735E]
const int INF = 1e9, MOD = 1e9 + 7;
vector<vector<vector<long long>>> dp;                 // dp[v][f][g], f up to k, g in [-1..k]
void ostap(int v, int p) {
    // init: v painted -> (f=0, g=-1); v not painted -> (f=INF, g=0).
    // for each child c:  combine  dp[v][f1][g1] * dp[c][f2][g2]  into ndp
    //   f_new = min(f1, f2 + 1);  g_new = max(g1, g2 + 1);
    //   if f_new + g_new <= k  then  g_new = -1  (all demands satisfied through v)
    // cap f_new and g_new by k.
}
// answer = sum of dp[root][f][g] over states with g == -1, mod 1e9+7.
\end{lstlisting}


\chaptersection{Write-From-Memory Drills}

Five signature skeletons, each lifted from a listing already in the
chapter. Cover the marked line, reproduce it from memory, then check the
answer one line below.

\begin{lstlisting}[caption={Drill 1 -- the capped child-merge (from \coderef{c27-knap})}]
void knap(int v, int p) {
    sz[v] = 1;
    dp[v].assign(V + 1, NEG);  dp[v][0] = /* FILL: seed */;
    for (int c : adj[v]) if (c != p) {
        knap(c, v);
        int lim  = /* FILL: min(sz[v], V) */;
        int clim = /* FILL: min(sz[c], V) */;
        vector<ll> ndp(V + 1, NEG);
        for (int i = 0; i <= lim; ++i) if (dp[v][i] > NEG)
            for (int j = 0; j <= clim && /* FILL: i + j <= V */; ++j)
                ndp[i + j] = max(ndp[i + j], dp[v][i] + dp[c][j]);
        dp[v].swap(ndp);
        sz[v] += sz[c];
    }
}
\end{lstlisting}
\noindent{\scriptsize Key: the two caps are the theorem. Drop either and
\thmref{c27-amortize} fails -- an uncapped star-root re-scans a fully-invalid
$\mathit{dp}$ array once per child, giving $\bigO{n^2}$ instead of $\bigO{nV}$.
The $i+j \le V$ guard is what keeps the merge inside the budget; forgetting
it is the most common WA on CF~855C (\probref{c27-knap-solved}).}

\begin{lstlisting}[caption={Drill 2 -- the rollback profile DFS (from \coderef{c27-profile})}]
int profile_dfs(int v, int p, int w) {
    int pos  = lower_bound(tail.begin(), tail.end(), w) - tail.begin();
    int back = /* FILL: (pos < tail.size() ? tail[pos] : -1) */;
    if (pos == (int)tail.size()) tail.push_back(w); else tail[pos] = w;
    int best = (int)tail.size();
    for (int u : adj[v]) if (u != p) best = max(best, profile_dfs(u, v, /* val */));
    /* FILL: undo the exact change */
      if (back == -1) tail.pop_back(); else tail[pos] = back;
    return best;
}
\end{lstlisting}
\noindent{\scriptsize Key: extend on entry, \emph{undo the exact slot on
exit}. Restoring the whole vector is $\bigO{\text{length}}$ and kills the
amortisation; overwriting only the touched position is $\bigO{1}$. Enumerate
every start vertex (or every peak) -- the profile is not a subtree
functional (\thmref{c27-profile-bound}, \probref{c27-profile-solved}).}

\begin{lstlisting}[caption={Drill 3 -- deepest-uncovered leaf-prune greedy (from \coderef{c27-greedy})}]
// order vertices by decreasing depth (BFS order reversed).
int answer = 0;
for (v in order_desc_by_depth) {
    if (/* FILL: v is a leaf AND dist_to_1(v) > R */) {
        int p = parent(v);
        attach_new_edge(p);                 // thm:c27-exchange: superset at p
        mark_covered_within_R(p);           // every node within R of p
        ++answer;
    }
}
\end{lstlisting}
\noindent{\scriptsize Key: attach at $\mathrm{parent}(v)$, not at $v$.
$\mathrm{cov}(\mathrm{parent}(v)) \supseteq \mathrm{cov}(v)$ is exactly the
containment that makes the greedy optimal (\thmref{c27-exchange}); swapping
$v$ and $\mathrm{parent}(v)$ is the top WA on CF~1029E. Root at~$1$ before
the depth sort.}

\begin{lstlisting}[caption={Drill 4 -- Grundy / Hackenbush fold (from \coderef{c27-grundy})}]
// generic mex
int grundy_mex(int v, int p) {
    vector<char> seen(/* FILL: adj[v].size() + 1 */);
    for (int u : adj[v]) if (u != p) seen[grundy_mex(u, v)] = 1;
    int g = 0; while (seen[g]) ++g; return g;
}
// Hackenbush closed form
int hackenbush(int v, int p) {
    int g = /* FILL: 0 */;                     // empty XOR at a leaf
    for (int u : adj[v]) if (u != p) g ^= /* FILL: (hackenbush(u,v) + 1) */;
    return g;
}
\end{lstlisting}
\noindent{\scriptsize Key: the $+1$ is the edge you stand on
(\thmref{c27-hackenbush}); the $\mathit{seen}$ array is sized
$\deg(v)+1$ because $g(v) \le \deg(v)$, not by a constant (Trap~2). Verdict:
first wins iff the root's value $\ne 0$ (\thmref{c27-grundy-xor}).}

\begin{lstlisting}[caption={Drill 5 -- two-flag trie DP (from \coderef{c27-455b-soln})}]
void fold_flags(int v) {
    if (trie[v].empty()) { win[v] = false; lose[v] = true; return; }
    win[v] = false; lose[v] = false;
    for (int u : trie[v]) {
        fold_flags(u);
        win[v]  |= /* FILL: !win[u] */;         // some child is a LOSS for opponent
        lose[v] |= /* FILL: !lose[u] */;        // some child is a WIN-FORCED-LOSS
    }
}
// verdict:  (!win[root]) -> "Second";
//           (lose[root]) -> "First";
//           else -> (k % 2 == 1 ? "First" : "Second").
\end{lstlisting}
\noindent{\scriptsize Key: leaf base is $(\mathit{win} = \mathit{false},
\mathit{lose} = \mathit{true})$ -- the mover cannot move and loses.
Both recurrences are OR-of-NOT (\thmref{c27-grundy-xor}'s Boolean shadow).
The $k$-parity case fires only when $\mathit{win}[\text{root}]$ is true
and $\mathit{lose}[\text{root}]$ is false -- First wins every game he
starts and cannot throw one (\probref{c27-game-solved}).}

\chaptersection{Complete Problem Checklist}

Every problem this chapter offers -- 4 fully solved in-section, 4
practice, one Summit. Tick \emph{only} after an unaided solve; a hinted
attempt counts as half. The ten mixed problems are listed in their own
station. Cite by name and number -- URLs appear once, in the solved and
practice entries.

\begin{center}
\footnotesize
\renewcommand{\arraystretch}{1.25}
\begin{tabular}{@{}c p{0.19\linewidth} l p{0.20\linewidth} p{0.235\linewidth} c@{}}
\toprule
\textbf{\#} & \textbf{Problem} & \textbf{Source} & \textbf{Key technique} & \textbf{What to look for} & \textbf{Done?} \\
\midrule
\multicolumn{6}{l}{\textbf{Section 14.1 -- tree knapsack, capped child-merge}} \\
W14.1.1 & Helga Hufflepuff's Cup & CF 855C & \lv{4} three-class $\bigO{n^2}$ tree knapsack & state (L/C/H) with budget $j$ = \# type-$k$ vaults; $\mathit{sz}$ caps the loops; \lstinline{long long} for the mod \\
P14.1.2 & Budgeted subtree (back-cite) & CF 161D & \lv{3} $\bigO{nk}$ histogram merge & Ch.13 owns the URL; re-solve cold with \thmref{c27-amortize}'s cap justification \\
\midrule
\multicolumn{6}{l}{\textbf{Section 14.2 -- states that grow, profile DFS}} \\
W14.2.1 & Treeland Tour & CF 490F & \lv{4} rollback patience LIS, per-start enumeration & state = tails array; extend on entry, undo exactly on exit; $\bigO{n^2 \log n}$ bound \\
P14.2.2 & (fresh value-indexed DP) & misc & \lv{4} growing profile & any DP whose state length is a parameter, not a coordinate \\
\midrule
\multicolumn{6}{l}{\textbf{Section 14.3 -- greedy on trees, exchange arguments}} \\
W14.3.1 & Tree with Small Distances & CF 1029E & \lv{4} deepest-uncovered leaf-prune greedy & attach at parent(v), not v; prove cov(parent(v)) superset cov(v) \\
P14.3.2 & Chloe and pleasant prizes & CF 743D & \lv{3} harvest two disjoint subtree sums & mx[v] = best subtree sum inside v; answer = best m1+m2 over a node's two child branches \\
\midrule
\multicolumn{6}{l}{\textbf{Section 14.4 -- games on trees}} \\
W14.4.1 & A Lot of Games & CF 455B & \lv{4} two-flag win/lose DP on a trie & leaf base (win=F, lose=T); verdict via $\mathit{lose}[\text{root}]$ or $k$-parity \\
P14.4.2 & Serval and Rooted Tree & CF 1153D & \lv{5} minimax as a sacrifice count & leaf $d=1$; max-node = min; min-node = sum; answer $= k + 1 - d[\text{root}]$ \\
\midrule
\multicolumn{6}{l}{\textbf{Summit (chapter end)}} \\
S14 & Ostap and Tree & CF 735E & \lv{5} supply/demand 2-index tree DP & state $(f,g)$ capped by $k$; reset $g=-1$ when $f+g \le k$; mod $10^9+7$ \\
\bottomrule
\end{tabular}
\end{center}

\chaptersection{Recommended Study Order}

Five focused days. Day~1 pays Ch.13's owed amortisation debt and internalises
the capped child-merge; day~2 lets the state grow into a profile and drills
rollback; day~3 exchanges a DP for a greedy and proves it; day~4 runs games
two ways (flags and Grundy); day~5 is the Summit. Solve in this order -- each
day's tool is the previous day's output.

\begin{center}
\small
\renewcommand{\arraystretch}{1.25}
\begin{tabular}{@{}llp{0.55\linewidth}@{}}
\toprule
\textbf{Day} & \textbf{Focus} & \textbf{Tasks} \\
\midrule
1 & Read + do 14.1 & trace the budget-$V{=}3$ knapsack on $T$ (\exref{c27-knap-run}); re-derive \thmref{c27-amortize} without looking at the proofbox; write \coderef{c27-knap} cold; CF~855C \\
2 & Read + do 14.2 & walk the rollback profile on $T$ (\exref{c27-profile-run}); CF~490F; verify the tails-array restore cost is $\bigO{1}$ per exit \\
3 & Read + do 14.3 & hand-simulate the deepest-uncovered greedy on $T$ with $R{=}2$ (\exref{c27-greedy-run}); CF~1029E; CF~743D as a two-branch harvest \\
4 & Read + do 14.4 & compute the Grundy overlay on $T$ (\exref{c27-grundy-run}); CF~455B; CF~1153D as a minimax-sacrifice count \\
5 & Summit + revise & CF~735E with the $(f,g)$ 2-index DP (\probref{c27-summit}); walk the decision flowchart on five new problem statements \\
\bottomrule
\end{tabular}
\end{center}

\chaptersection{Difficulty Progression}

Tiers are competency gates -- each row names the \emph{skill} the tier
unlocks, not just a rating bucket.

\begin{center}
\small
\renewcommand{\arraystretch}{1.3}
\begin{tabular}{@{}lp{0.36\linewidth}p{0.34\linewidth}@{}}
\toprule
\textbf{Tier} & \textbf{Problems} & \textbf{What passing this tier proves} \\
\midrule
$\lv{3}$ \ $\bigstar\bigstar\bigstar$
  & CF~743D, CF~1914F, CF~274B
  & you can pick the right DP shape among the four and hand-code the state without a template \\
$\lv{4}$ \ $\bigstar\bigstar\bigstar\bigstar$
  & CF~855C, CF~490F, CF~1029E, CF~455B, CF~1153D, CF~915F, CF~588E, CF~696A, CF~1065F, CF~2063E, CF~768E
  & you cap the merge loops, roll back a profile exactly, prove an exchange by a superset containment, and choose flags or Grundy on a game \\
$\lv{5}$ \ $\bigstar\bigstar\bigstar\bigstar\bigstar$
  & CF~735E (Summit), CF~543D, CF~613D
  & you compose a two-index supply/demand state and reason about the reset condition; you \emph{recognise} when a problem is really a Ch.15 rerooting or Ch.17 DSU-on-tree in disguise \\
\bottomrule
\end{tabular}
\end{center}

\chaptersection{Reduction Gallery}

Every DP here is ``fix a root, name a state, fold over children -- but pay for
the state's size.'' Each line collapses a family onto a listing already in
the chapter -- no new theory.

\begin{itemize}[nosep, leftmargin=1.4em]
  \item ``colour the tree so that each subtree has consecutive values, count
        assignments'' (CF~855C) $\longrightarrow$ three-class knapsack with a
        type-$k$-count budget; the merge is \coderef{c27-knap} with a
        multinomial coefficient (\probref{c27-knap-solved}).
  \item ``longest increasing path in a tree'' (CF~490F) $\longrightarrow$
        rollback patience profile with per-start enumeration --
        \thmref{c27-profile-bound}'s $\bigO{n^2 \log n}$ bound
        (\probref{c27-profile-solved}).
  \item ``add minimum edges so every leaf is within $R$ of the root''
        (CF~1029E) $\longrightarrow$ deepest-uncovered leaf-prune greedy,
        proved by the superset containment of \thmref{c27-exchange}
        (\probref{c27-greedy-solved}).
  \item ``two players pick ropes and take subtrees; maximise the sum of both''
        (CF~743D) $\longrightarrow$ subtree-sum $\mathit{mx}[v]$ + best two
        disjoint child branches at some $v$ -- a peak-harvest instance of
        Ch.13's \propref{c10-pairtop} at a different $\Phi$.
  \item ``trie game, $k$ rounds chained by loser-starts-next'' (CF~455B)
        $\longrightarrow$ two-flag win/lose DP plus a parity case split
        (\probref{c27-game-solved}).
  \item ``minimax tree where internal nodes are $\max$ or $\min$'' (CF~1153D)
        $\longrightarrow$ sacrifice-count DP: leaf $d=1$, max-node min, min-node
        sum, answer $k + 1 - d[\text{root}]$ (\secref{c27-games} practice).
  \item ``paint $\le k$-spread black seeds so every vertex is covered; count
        colourings'' (CF~735E) $\longrightarrow$ supply/demand 2-index
        tree DP with the $f+g \le k$ reset (\probref{c27-summit}).
\end{itemize}

\chaptersection{Mixed Problem Set}

These refuse to say ``DP''; decide \emph{which quantity lives at each node
and what $\Phi$ is} before opening the code. A hint tells you the shape of
the answer, never the tool.

\begin{mixedset}
  \practice{Zero Tree}{Codeforces 274B}
    {https://codeforces.com/problemset/problem/274/B}{CF 1700 · L4}
    {A tree with integer values; each move adds $+1$ or $-1$ to every node in a chosen subtree. Minimum moves to make every value zero. State per node = (increment, decrement) counts required by the subtree; children fold by $\max$ per coordinate -- the merge of \secref{c27-knapsack} with the $\max$-combiner replacing the $+$-convolution (\coderef{c27-kit} row~1 variant).}
  \practice{Road Decomposition}{Codeforces 543D}
    {https://codeforces.com/problemset/problem/543/D}{CF 2300 · L5}
    {For each query root $x$, count ways to designate edges ``good'' so every path from $x$ uses at most one bad edge. Per-root DP is one knapsack-shaped fold; the multi-root sweep is Ch.15's rerooting. Do the single-root version here -- budget the state by child index and merge one edge at a time.}
  \practice{Duff in the Army}{Codeforces 588E}
    {https://codeforces.com/problemset/problem/588/E}{CF 2100 · L4}
    {Multiple queries: along a path $u \to v$, list the top-$10$ minimum general IDs on that path. Cap the profile state at~$10$ and merge along the path -- an instance of \secref{c27-growing}'s ``grow the state to a bounded profile'' with LCA as a Ch.16 forward-pointer.}
  \practice{Imbalance Value of a Tree}{Codeforces 915F}
    {https://codeforces.com/problemset/problem/915/F}{CF 2100 · L4}
    {Sum over all simple paths of (max edge $-$ min edge) on the path. Linearity of expectation splits it into $\sum$ max $-$ $\sum$ min; each half is a Kruskal-style union-find sweep on the tree's edges sorted by weight -- a profile DP indexed by threshold. Compare with Ch.13's peak-harvest geometry.}
  \practice{Lorenzo Von Matterhorn}{Codeforces 696A}
    {https://codeforces.com/problemset/problem/696/A}{CF 1800 · L4}
    {Infinite heap-tree (parent = $u/2$): path updates $u \leftrightarrow v$ by $\pm 1$, path sum queries. Rooting at~$1$ collapses to LCA arithmetic -- Ch.16's \chref{c11}; the per-node state is a scalar $+$ a subtree sum, exactly \thmref{c10-subtree-template} at $V=1$, but the update rule is the ``growing profile'' of \secref{c27-growing} in reverse.}
  \practice{Kingdom and its Cities}{Codeforces 613D}
    {https://codeforces.com/problemset/problem/613/D}{CF 2200 · L5}
    {Minimum non-important vertices to remove so that no two important vertices remain connected. The greedy is deepest-first (Ch.14's \thmref{c27-exchange}) with a per-node ``is there an unblocked important below?'' flag; the DP pair (has-below, has-above) matches the two-state style of Ch.13's \secref{c10-twostate}.}
  \practice{Up and Down the Tree}{Codeforces 1065F}
    {https://codeforces.com/problemset/problem/1065/F}{CF 2100 · L5}
    {A token moves down to any leaf or up $\le k$ edges (never twice in a row). Maximise distinct leaves visited. The state at $v$ is a profile indexed by how many ``up-moves'' remain -- \secref{c27-growing} with the budget cap $k$ of \thmref{c27-amortize}.}
  \practice{Programming Competition}{Codeforces 1914F}
    {https://codeforces.com/problemset/problem/1914/F}{CF 1400 · L3}
    {The problem the plan originally slotted into 14.3; live study showed the greedy is a path-invariant hybrid rather than a clean superset exchange, so it lives here instead. Root the tree, count ``good'' leaves per subtree by a scalar fold -- \thmref{c27-amortize}'s merge with $V=1$, i.e.~Ch.13's shape.}
  \practice{Triangle Tree}{Codeforces 2063E}
    {https://codeforces.com/problemset/problem/2063/E}{CF 1900 · L4}
    {Sum $f(u,v)$ over ``good'' (non-ancestor) pairs $(u,v)$; $f$ is a combinatorial function of the depths. Peak-harvest identity (\propref{c10-pairtop}) with $F(u,v)$ = the pair's contribution; the state at $\ell$ is a depth histogram, so the merge is \secref{c27-knapsack}'s convolution capped by $\height$.}
  \practice{Game of Stones}{Codeforces 768E}
    {https://codeforces.com/problemset/problem/768/E}{CF 1800 · L4}
    {$n$ piles, each a Nim-variant where a move takes $1$ or $2$ stones but never the same count twice in a row. Each pile's Grundy value $s \le 60$ fits a small $\operatorname{mex}$ DP; XOR across piles (thm:c27-grundy-xor). The tree is a star at the ``virtual ground''; the fold is \coderef{c27-grundy} instantiated on independent heaps.}
\end{mixedset}


\chaptersection{The Summit -- CF 735E Ostap and Tree}

\begin{notebox}[Why this is the chapter top rope]
Sections 14.1--14.4 each carry \emph{one} growing index: a budget, a
profile length, a coverage radius, a Grundy value. The Summit carries
\emph{two} simultaneously -- a \textbf{supply} $f$ (distance from $v$ to
the nearest black node inside $v$'s subtree, $\infty$ if none) and a
\textbf{demand} $g$ (distance from $v$ to the farthest still-uncovered
node inside $v$'s subtree, $-\infty$ if none). Every subtree DP that Ch.13
and Ch.14 have taught so far is the special case where one of $f, g$ is
absent. The reset rule -- whenever $f + g \le k$, the demand side collapses
to ``nothing outstanding'' -- is the pair-local version of
\thmref{c27-amortize}'s budget cap, applied to both coordinates at once,
and it is what makes the whole table $O(nk^2)$ instead of $O(n^3)$. Every
tool in the answer is already in \coderef{c27-kit}: the capped child-merge
of \secref{c27-knapsack}, the reset-condition from \thmref{c27-exchange}'s
superset argument, and the mod-$10^9{+}7$ counting fold of \secref{c27-growing}.
\end{notebox}

\solvedproblem{c27-summit}{Ostap and Tree}{Codeforces 735E}
  {https://codeforces.com/problemset/problem/735/E}{CF 2400 · L5}
  {two-index supply/demand tree knapsack, capped-by-k reset rule, mod 1e9+7 counting fold, deepest-cover exchange in a 2-D state}

\begin{problemstatement}[Ostap and Tree (statement and samples verified against the live problem)]
A tree of $n$ vertices. Choose a subset $B \subseteq V$ to paint black so
that every vertex lies within distance $\le k$ of some element of $B$
(equivalently: $B$ is a distance-$k$ dominating set). Count such subsets
modulo $10^9+7$. Constraints: $1 \le n \le 100$; $0 \le k \le
\min(20, n-1)$.\\
\textbf{Sample~1}: $n = 2$, $k = 0$, edge $1$--$2$ $\Rightarrow$ output
\texttt{1}: only $\{1,2\}$ works, since a single vertex's distance-$0$
neighbourhood is itself. \textbf{Sample~2}: $n = 2$, $k = 1$, edge
$1$--$2$ $\Rightarrow$ output \texttt{3}: $\{1\}$, $\{2\}$ and $\{1,2\}$
all cover (each vertex has a black at distance $\le 1$); $\varnothing$
does not. \textbf{Sample~3}: $n = 4$, $k = 1$, path $1$--$2$--$3$--$4$
$\Rightarrow$ output \texttt{9}: the $9$ dominating sets are
$\{1,3\}, \{1,4\}, \{2,3\}, \{2,4\}, \{1,2,3\}, \{1,2,4\}, \{1,3,4\},
\{2,3,4\}, \{1,2,3,4\}$.
\end{problemstatement}

\begin{attemptgate}[Attempt it yourself first -- \lv{5}, give yourself 60 minutes]
Naive: enumerate subsets, $O(2^n \cdot n)$ -- dead at $n = 100$. Root at~$1$
and try a subtree DP. A single-index state cannot encode both ``is there a
black below? how far?'' and ``is there a still-uncovered node below? how
far?'' \emph{simultaneously}. Two questions unlock the shape: (i) what two
distances must $v$'s state report to its parent so the parent can decide,
without knowing $v$'s subtree structure? (ii) when is a demand
``absorbed'' inside $v$'s subtree?
\peek{If the nearest black below $v$ is at distance $f$ and the
      farthest uncovered-below is at distance $g$, when can $v$ report
      $g = -\infty$ (nothing outstanding) to its parent?}
\nudge{Exactly when the black at distance $f$ also covers that farthest
      uncovered node, i.e.\ when $f + g \le k$. The reset is what collapses
      the two-index state to a single value at a node that "seals" its
      subtree.}
\fullhint{State $\mathit{dp}[v][f][g]$: number of colourings of $v$'s
      subtree in which the nearest black to $v$ inside the subtree is at
      distance $f$ (capped at $k$; use $f = \infty$ if none) and the
      farthest still-uncovered node inside the subtree is at distance $g$
      from $v$ (capped at $k$; $g = -1$ sentinel for ``none''). Leaf:
      paint black gives state $(0, -1)$; paint white gives $(\infty, 0)$.
      Merge a child $c$ by combining $(f_1, g_1)$ at the running $v$
      with $(f_2, g_2)$ at $c$: $f' = \min(f_1, f_2 + 1)$,
      $g' = \max(g_1, g_2 + 1)$; if $f' + g' \le k$ then reset $g' = -1$.
      Prune any state with $g' > k$ (such a demand can no longer be met
      through the parent). Multiply by $2$ ways of painting $v$ at the
      leaf; sum-product over combinations; mod $10^9 + 7$. Answer
      $\sum_{f} \mathit{dp}[\text{root}][f][-1]$. Time $O(nk^2)$ by
      \thmref{c27-amortize}'s budget cap on both coordinates. Full build
      below.}
\end{attemptgate}

\solstep{Understanding the Problem}
The predicate is a covering constraint coupling $v$ to every descendant and
every ancestor simultaneously: a white vertex is legal exactly when \emph{some}
black node is within distance $k$. From the perspective of $v$'s subtree, two
things matter -- \emph{how near is a supply of blackness (either at $v$ or below
$v$), and how far reaches an unmet demand for coverage from below?} A state
that keeps both distances (bounded by $k$) can be handed up to $v$'s parent
with no other information: the parent only needs to know how far up the supply
and demand rays extend. If they cross -- $f + g \le k$ -- the demand is met
inside $v$'s subtree and can be reported as ``none outstanding'' to the parent.

\solstep{Key Observations}
\begin{observation}
Two coordinates suffice per node: $f = \min\{$distance from $v$ to a black node
in $v$'s subtree$\}$ (capped at $k$; $\infty$ if none in the subtree) and
$g = \max\{$distance from $v$ to a still-uncovered node in $v$'s subtree$\}$
($-1$ if none; capped at $k$).
\end{observation}
\begin{observation}
A demand $g$ that survives past a node $v$ to its parent must satisfy $g \le k$;
if it grows to $g > k$ at $v$ (meaning an uncovered node is farther than $k$ from
every black in the current view), the branch cannot recover from above within
distance $k$ either -- such states are \emph{pruned}.
\end{observation}
\begin{observation}
Merge rule at $v$ with running state $(f_1, g_1)$ and child state $(f_2, g_2)$
attached by an edge: $f' = \min(f_1, f_2 + 1)$, $g' = \max(g_1, g_2 + 1)$; then
if $f' + g' \le k$, reset $g' = -1$ (the supply reaches the demand through $v$).
The $+1$ is the edge you cross -- the same $+1$ as \thmref{c27-hackenbush}, now
attached to a coverage rather than a nimber.
\end{observation}
\begin{observation}
Base at a leaf: painting the leaf black gives $(0, -1)$; painting it white gives
$(\infty, 0)$ -- an uncovered node at distance $0$ from itself.
\end{observation}
\begin{observation}
Answer at the root: every demand must be sealed inside the tree, so sum
$\mathit{dp}[\text{root}][f][-1]$ over all supply values $f \in \{0, \dots, k, \infty\}$.
\end{observation}

\solstep{Think It Through With Me}
\begin{thinkbox}
\deadend{Standard Ch.13 subtree DP with a single colour-per-node state, plus
an ``is covered?'' flag propagated from the parent.}
{The flag cannot tell the parent \emph{how much distance remains} for a supply
to reach down and cover an unmet demand in a sibling subtree. You need two
integers per node, both capped by $k$.}
\deadend{Enumerate a set $B$ and check via BFS from every $b \in B$.}
{$\bigO{2^n \cdot n}$; at $n = 100$ this is $\sim 10^{32}$ -- hopeless.}
\deadend{Greedy: paint the deepest node, propagate coverage, recurse.}
{Fails Sample~3 -- greedy paints $\{2, 4\}$ but not necessarily all $9$; you
must \emph{count} dominating sets, not find a minimum one. The Summit is a
\probref{c27-knap-solved} flavour (counting fold), not a
\probref{c27-greedy-solved} flavour (extremal greedy).}
\aha{Two-index state $(f,g)$ capped by $k$; reset $g$ to $-1$ whenever the
supply ray crosses the demand ray at the current node; one capped child-merge
per edge in the shape of \thmref{c27-amortize}. Answer at $g = -1$ at the root.}
\end{thinkbox}

\solstep{The Approach Ladder}
\begin{center}\renewcommand{\arraystretch}{1.3}
\begin{tabular}{@{}c p{0.235\linewidth} p{0.30\linewidth} l c@{}}
\toprule
\textbf{\#} & \textbf{Approach} & \textbf{Idea} & \textbf{Time} & \textbf{Verdict} \\
\midrule
1 & \lv{3} Subset enumeration & test each $B \subseteq V$ by BFS & $O(2^n n)$ & \tle \\
2 & \lv{4} Single-index colour DP & colour $v$ black/white, one flag ``covered?'' & $O(n)$ & \textcolor{warningred}{\ding{55}\ \textbf{WA}} \\
3 & \lv{5} Two-index supply/demand DP & $(f, g)$ capped by $k$, reset on $f + g \le k$ & $O(n k^2)$ & \ac \\
\bottomrule
\end{tabular}\end{center}
Rung~1 is honest but exponential. Rung~2 misses the distance information
needed to decide coverage across sibling subtrees -- the flag is a
Ch.13-style constant tuple in a regime where \thmref{c27-amortize} demands an
index. Rung~3 is \emph{this} chapter's answer -- \secref{c27-knapsack}'s
capped merge, doubled to two coordinates, with \thmref{c27-exchange}'s
superset-containment idea promoted from a greedy proof to a state reset.

\solstep{Solution}
Root at~$1$. Post-order DFS. At node $v$: build $\mathit{dp}[v]$ as a table of
size $(k+2) \times (k+1)$ (rows: supply $f \in \{0, \dots, k, \infty\}$ encoded
as index $k+1$; columns: demand $g \in \{-1, 0, \dots, k\}$). Initialise at $v$:
two entries $(0, -1) \mathrel{+}= 1$ (paint $v$ black) and
$(\infty, 0) \mathrel{+}= 1$ (paint $v$ white). For each child $c$, first DFS
into $c$; then combine running $\mathit{dp}[v][f_1][g_1]$ with $\mathit{dp}[c][f_2][g_2]$:
$f' = \min(f_1, f_2 + 1)$, $g' = \max(g_1, g_2 + 1)$; if $f' + g' \le k$
set $g' = -1$; skip if $g' > k$ or $f' > k$; add
$\mathit{dp}[v][f_1][g_1] \cdot \mathit{dp}[c][f_2][g_2]$ into
$\mathit{ndp}[f'][g']$. Swap $\mathit{ndp}$ into $\mathit{dp}[v]$. Answer:
$\sum_{f} \mathit{dp}[1][f][-1] \bmod 10^9{+}7$. Cap every loop bound by
$k$; this is \thmref{c27-amortize} in two dimensions -- $\bigO{n k^2}$ at
$n \le 100$, $k \le 20$ is $4 \cdot 10^4$ operations.

\solstep{Proof of Correctness}
\begin{proofbox}[\coderef{c27-summit} counts exactly the distance-$k$ dominating sets]
\emph{Soundness of the state.} Fix $v$. The pair $(f,g)$ is defined so that the
parent of $v$ can complete the coverage test \emph{without knowing $v$'s subtree
shape}: a black at distance $f$ from $v$ covers every vertex at distance $\le k - f$
above $v$'s subtree; a demand at distance $g$ from $v$ must be met at distance
$\le k - g$ above $v$'s subtree. Thus only the two scalars $f, g$ are needed, and
no two colourings with different $(f,g)$ can be re-unified above $v$. \emph{Merge
correctness.} Consider a partial colouring of $v$'s processed subtree (state $(f_1, g_1)$)
and a colouring of $c$'s subtree (state $(f_2, g_2)$). After merging $c$ under $v$,
every distance from $v$ to a supply in $c$'s subtree is $f_2 + 1$ (crossing the edge
$(v,c)$), and to a demand in $c$'s subtree is $g_2 + 1$; the running $v$-side
supplies and demands are unchanged. So $f' = \min(f_1, f_2 + 1)$, $g' = \max(g_1, g_2 + 1)$.
The \emph{reset} $g' \gets -1$ fires exactly when every demand in the merged subtree is
within $k$ of some supply in it -- i.e., the farthest demand $g'$ is covered by the
nearest supply $f'$ through $v$'s node itself, since the shortest path from a supply
at distance $f'$ to a demand at distance $g'$ passing through $v$ has length $f' + g'$.
If $f' + g' \le k$, the demand is met internally; if $f' + g' > k$, the demand still
needs a supply from above. \emph{Pruning.} If $g' > k$ after the merge, no supply
outside the subtree can reach the farthest demand either (the parent's view adds at
least $1$ more distance per level above $v$, and the subtree is already $> k$) --
so this branch is not a valid dominating set at any completion and can be dropped.
\emph{Leaf base.} Painting leaf $v$ black: $f = 0$ (the black at $v$ itself), $g = -1$
(no demand uncovered by $v$). Painting $v$ white: no supply inside $v$'s subtree, so
$f = \infty$; the vertex $v$ itself is uncovered at distance $0$, so $g = 0$. \emph{Answer.}
At the root, every demand must be met \emph{inside} the whole tree (there is no ``above''
to supply coverage), so we sum exactly the states with $g = -1$. \qed
\end{proofbox}

\solstep{Implementation}
\begin{lstlisting}[caption={Ostap and Tree -- supply/demand two-index tree DP, capped by $k$, with the $f+g \le k$ reset; count mod $10^9{+}7$ (Codeforces 735E; $O(n k^2)$)}, label={lst:c27-summit}]
#include <bits/stdc++.h>
using namespace std;
typedef long long ll;
const ll MOD = 1000000007LL;
int n, K;
vector<vector<int>> adj;

// dp[v][f][g]: #ways in subtree(v) with supply f (0..K, K+1 = INF) and
//              demand g (-1 = NONE, 0..K).  Array g shifted by +1.
// size (K+2) x (K+2).
vector<vector<vector<ll>>> dp;

void dfs(int v, int p) {
    dp[v].assign(K + 2, vector<ll>(K + 2, 0));
    // Base: paint v black -> (f=0, g=-1 encoded as index 0).
    dp[v][0][0] = 1;
    // Paint v white -> (f=INF = K+1, g=0 encoded as index 1).
    dp[v][K + 1][1] = 1;
    for (int c : adj[v]) if (c != p) {
        dfs(c, v);
        vector<vector<ll>> ndp(K + 2, vector<ll>(K + 2, 0));
        for (int f1 = 0; f1 <= K + 1; ++f1)
            for (int g1 = 0; g1 <= K + 1; ++g1) if (dp[v][f1][g1]) {
                int realg1 = g1 - 1;                // -1 or 0..K
                for (int f2 = 0; f2 <= K + 1; ++f2)
                    for (int g2 = 0; g2 <= K + 1; ++g2) if (dp[c][f2][g2]) {
                        int realg2 = g2 - 1;
                        int fnew = min(f1, f2 == K + 1 ? K + 1 : min(K + 1, f2 + 1));
                        int gnew;
                        if (realg1 < 0 && realg2 < 0) gnew = -1;
                        else gnew = max(realg1, realg2 == K + 1 ? K + 1 : realg2 + 1);
                        // (Note: realg2 capped by K, so realg2+1 is only
                        //  used when a child demand is still open.)
                        // Actually demands from child arrive shifted by +1.
                        int dchild = (realg2 < 0) ? -1 : realg2 + 1;
                        gnew = max(realg1, dchild);
                        if (gnew > K) continue;                        // prune
                        if (fnew <= K && gnew >= 0 && fnew + gnew <= K)
                            gnew = -1;                                // reset
                        int gi = gnew + 1;                             // re-encode
                        ndp[fnew][gi] = (ndp[fnew][gi]
                            + dp[v][f1][g1] * dp[c][f2][g2]) % MOD;
                    }
            }
        dp[v].swap(ndp);
    }
}

int main(){
    if (scanf("%d %d", &n, &K) != 2) return 0;
    adj.assign(n + 1, {});
    dp.assign(n + 1, {});
    for (int i = 0; i < n - 1; ++i){
        int a, b; scanf("%d %d", &a, &b);
        adj[a].push_back(b); adj[b].push_back(a);
    }
    dfs(1, 0);
    ll ans = 0;
    for (int f = 0; f <= K + 1; ++f) ans = (ans + dp[1][f][0]) % MOD;   // g=-1 -> idx 0
    printf("%lld\n", ans);
    return 0;
}
\end{lstlisting}

\begin{notebox}[The four lines that are the whole idea]
\code{int fnew = min(f1, f2 + 1);} -- the supply ray crosses the edge.
\code{int gnew = max(realg1, realg2 + 1);} -- the demand ray crosses the edge.
\code{if (fnew + gnew <= K) gnew = -1;} -- the rays meet at $v$.
\code{if (gnew > K) continue;} -- the demand outran any supply from above.
Everything else is array plumbing. This is \thmref{c27-amortize} doubled:
\emph{both} coordinates capped by $k$, and the merge is one child at a time.
\end{notebox}

\solstep{Trace on Sample 1}
\begin{tracebox}[Sample~1 -- $n = 2$, $k = 0$, edge $1$--$2$ (answer $1$)]
Root at~$1$; $1$ is parent of~$2$; $K=0$. Size table: $(K+2) \times (K+2) = 2 \times 2$,
indices $f \in \{0, \infty\}$, $g \in \{-1, 0\}$ encoded $\{0, 1\}$.

\textbf{DFS at leaf~$2$}: $\mathit{dp}[2][0][0] = 1$ (black, state $(0, -1)$);
$\mathit{dp}[2][\infty][1] = 1$ (white, state $(\infty, 0)$).

\textbf{DFS at $1$}: init $\mathit{dp}[1][0][0] = 1$ (black at~$1$), $\mathit{dp}[1][\infty][1] = 1$ (white at~$1$).
Merge child~$2$ (running $v$-state $\times$ child state):

\emph{Case A}: $v$ black, $f_1 = 0$, $g_1 = -1$. Cross with child.
  \begin{itemize}[nosep]
    \item[(a1)] child black, $(f_2 = 0, g_2 = -1)$: $f' = \min(0, 0+1) = 0$;
      $g' = \max(-1, -1) = -1$; not pruned; reset not needed. Encode: $\mathit{ndp}[0][0] \mathrel{+}= 1$.
    \item[(a2)] child white, $(f_2 = \infty, g_2 = 0)$: $f' = \min(0, \infty + 1) = 0$;
      $g' = \max(-1, 0 + 1) = 1$. Prune since $g' = 1 > K = 0$. Discard.
  \end{itemize}
\emph{Case B}: $v$ white, $f_1 = \infty$, $g_1 = 0$. Cross with child.
  \begin{itemize}[nosep]
    \item[(b1)] child black, $(f_2 = 0, g_2 = -1)$: $f' = \min(\infty, 0+1) = 1$.
      Cap by $K+1$: but $f' = 1 > K = 0$, so we mark it as $\infty$ (supply too far to help);
      keep $f' = K+1 = 1$ (index $\infty$). $g' = \max(0, -1) = 0$. Reset: $f' + g' = 1 + 0 = 1 > K$, no reset.
      Encode: $\mathit{ndp}[1][1] \mathrel{+}= 1$ -- state $(\infty, 0)$.
    \item[(b2)] child white, $(f_2 = \infty, g_2 = 0)$: $f' = \infty$; $g' = \max(0, 1) = 1$.
      Prune since $g' = 1 > K = 0$. Discard.
  \end{itemize}

Result at root: $\mathit{ndp}[0][0] = 1$ (both black), $\mathit{ndp}[1][1] = 1$ (only~$2$ black, but $v = 1$ is uncovered).
Answer: $\sum_f \mathit{ndp}[f][0]$ (only $g = -1$ counts as valid) $= \mathit{ndp}[0][0] = 1$. \checkmark
\end{tracebox}

\solstep{Trace on Sample 2}
\begin{tracebox}[Sample~2 -- $n = 2$, $k = 1$, edge $1$--$2$ (answer $3$)]
$K = 1$; table size $3 \times 3$, indices $f \in \{0, 1, \infty\}$, $g \in \{-1, 0, 1\}$.

\textbf{DFS at leaf~$2$}: $\mathit{dp}[2][0][0] = 1$ (black, state $(0, -1)$);
$\mathit{dp}[2][2][1] = 1$ (white, state $(\infty, 0)$, encoded $f = 2 = K+1$, $g = 1 = 0{+}1$).

\textbf{DFS at $1$}: init $\mathit{dp}[1][0][0] = 1$, $\mathit{dp}[1][2][1] = 1$.

\emph{Merge child~$2$ into $\mathit{dp}[1][0][0]$ (running $v$ black)}:
  \begin{itemize}[nosep]
    \item child black: $f' = 0$; $g' = \max(-1, -1) = -1$; no reset. Encode $\mathit{ndp}[0][0] \mathrel{+}= 1$. $\to \{1, 2\}$ both black.
    \item child white: $f' = \min(0, \infty+1) = 0$; $g' = \max(-1, 0+1) = 1$; reset test $f' + g' = 1 \le K = 1$: \textbf{reset} $g' = -1$. Encode $\mathit{ndp}[0][0] \mathrel{+}= 1$. $\to \{1\}$, valid.
  \end{itemize}

\emph{Merge child~$2$ into $\mathit{dp}[1][2][1]$ (running $v$ white, $f_1 = \infty$, $g_1 = 0$)}:
  \begin{itemize}[nosep]
    \item child black: $f' = \min(\infty, 0+1) = 1$; $g' = \max(0, -1) = 0$; reset $1 + 0 = 1 \le 1$: \textbf{reset} $g' = -1$. Encode $\mathit{ndp}[1][0] \mathrel{+}= 1$. $\to \{2\}$, valid.
    \item child white: $f' = \min(\infty, \infty) = \infty$; $g' = \max(0, 1) = 1$; no supply ($f' > K$) so no reset. Encode $\mathit{ndp}[2][2] \mathrel{+}= 1$ -- invalid at root ($g \ne -1$).
  \end{itemize}

Answer: $\sum_f \mathit{ndp}[f][0] = 1 + 1 + 1 = \mathbf{3}$. \checkmark
\end{tracebox}

\solstep{Trace on Sample 3 (sketch)}
\begin{tracebox}[Sample~3 -- $n = 4$, $k = 1$, path $1$--$2$--$3$--$4$ (answer $9$)]
The full DP is $O(n k^2) = 4 \cdot 4 = 16$ state entries; running the merge
line-by-line is mechanical and not printed here. Instead, \emph{cross-check by
independent enumeration}: a dominating set of a path $P_4$ with $k = 1$ must
place blacks so that every vertex is within distance~$1$ of one.
Non-trivial minimal dominating sets: $\{1,3\}, \{1,4\}, \{2,3\}, \{2,4\}$.
Each extends to a set by adding any superset vertex: $\{1,2,3\}, \{1,3,4\},
\{1,2,4\}, \{2,3,4\}, \{1,2,3,4\}$. That is $4 + 5 = 9$ dominating sets, which
matches the published sample output $9$. \checkmark
\end{tracebox}

\solstep{Complexity}
\begin{complexitybox}
Time $\bigO{n k^2}$: for each node $v$ and each child $c$, the merge visits every
$(f_1, g_1)$ paired with every $(f_2, g_2)$, four nested loops over $[0, k] \cup \{\infty\}$.
By \thmref{c27-amortize}'s argument the total over the tree is $\bigO{n k^2}$ --
no vertex visits its state more than $\bigO{\text{degree}}$ times, and each state slot
touches each child slot exactly once. Memory $\bigO{n k^2}$ for the $\mathit{dp}$ tables
(or $\bigO{k^2}$ if you free on exit). At $n = 100$, $k \le 20$: $4 \cdot 10^4$ operations
-- trivially fast. Verdict: \ac.
\end{complexitybox}

\solstep{Verification and Edge Cases}
\begin{center}\renewcommand{\arraystretch}{1.2}
\begin{tabular}{@{}p{0.30\linewidth} p{0.62\linewidth}@{}}
\toprule
\textbf{Case} & \textbf{What the code does} \\
\midrule
$k = 0$ & Every vertex must be black to be covered by itself; the DP collapses to a single valid state (all black) at the root, giving answer~$1$ for any connected tree. Matches Sample~1. \\
$k \ge$ diameter & Every subset that contains at least one black is a dominating set; answer is $2^n - 1$. Check on a star with $k = 1$: paint the centre or any leaf, or both. \\
$n = 1$ & The single vertex is its own subtree; leaf base fires twice, states $(0, -1)$ and $(\infty, 0)$; sum over $g = -1$ at the root is $1$. \\
Chain of length~$n$, $k = 1$ & Every other vertex must be painted or one adjacent; the count is a Fibonacci-shaped recurrence in $n$; the DP reproduces it. \\
Overflow & $\mathit{dp}$ entries are $\bmod 10^9{+}7$ at every add; the product of two entries fits in \lstinline{long long} ($< (10^9)^2 < 10^{18}$). \\
\bottomrule
\end{tabular}\end{center}

\solstep{Pitfalls and Variants}
\begin{warningbox}
Ways to break it: (i) \emph{missing the $+1$ shift} on the child's supply or demand
-- the merge must add~$1$ to $f_2$ and to $g_2$ (the edge $(v, c)$ is real distance).
(ii) \emph{reset tested before the shift} -- apply $\min(f_1, f_2+1)$ and
$\max(g_1, g_2+1)$ \emph{first}, then test $f' + g' \le k$. (iii) \emph{not pruning}
$g' > k$ -- states with an unreachable demand explode the count and break mod
arithmetic on the sample. (iv) \emph{index confusion on $g = -1$} -- the array
column is $g+1$ so index~$0$ means $-1$; the answer sums over column~$0$ at the root.
\emph{Variants:} the ``minimum number of black vertices'' version replaces the counting
fold by a minimisation over $f$ and $g$ and drops the mod; the ``count colourings with
exactly $b$ blacks'' version adds a third coordinate to the state and pays
$\bigO{n b k^2}$; and the general ``distance-$k$ set cover on a tree'' family is a
recurring Ch.15 (\chref{c20}) pattern once the coverage must hold for every root.
\end{warningbox}

\begin{insightbox}[Key takeaways -- three things to remember]
\begin{enumerate}[nosep, leftmargin=1.6em]
  \item \emph{Two rays per node: supply and demand.} Whenever a tree DP
        must \emph{both} witness coverage from below and remember what still
        needs coverage from above, the state doubles to $(f, g)$; the reset
        condition $f + g \le k$ is the pair-local version of
        \thmref{c27-exchange}'s superset containment (\probref{c27-summit}).
  \item \emph{Capping both coordinates is the amortisation.}
        \thmref{c27-amortize} applies per index; two indices, both capped by~$k$,
        give $\bigO{n k^2}$ total. Forgetting to cap $g$ at $k$ is a
        correctness bug, not just a performance one -- an un-capped demand
        cannot be met from above and must be pruned.
  \item \emph{Counting folds look like optimisation folds.} Replace $\max$
        by $+$, $\mathit{NEG\_INF}$ seed by~$0$, and the answer extraction
        by a sum over legal final states -- the same \coderef{c27-knap}
        shape that solved CF~855C's counting problem in \secref{c27-knapsack}
        is what the Summit re-rolls in a doubled state.
\end{enumerate}
\end{insightbox}

\begin{hardeningbox}
\begin{itemize}[nosep, leftmargin=1.4em]
  \item \textbf{Generalises}: any ``distance-$k$ cover on a tree, count or
        optimise'' collapses to a supply/demand pair capped by~$k$. Ch.15's
        rerooting replays the trick across all roots; Ch.24's centroid
        decomposition runs it across path-length buckets (\chref{c20}, \chref{c18}).
  \item \textbf{The six listings of this chapter are one skeleton}:
        \coderef{c27-knap}, \coderef{c27-855c-soln}, \coderef{c27-profile},
        \coderef{c27-greedy}, \coderef{c27-grundy}, \coderef{c27-455b-soln}
        all share ``one DFS, one per-frame state, one merge -- capped by a
        parameter''; the Summit is a seventh instantiation with a two-coordinate
        state and a reset (see \coderef{c27-kit}).
  \item \textbf{Alternatives}: a $2$-SAT encoding of the coverage predicate
        gives $O(n \cdot 2^k)$ per satisfying assignment count -- unusable
        at $k \le 20, n \le 100$ in a 1-second limit; the DP is the honest
        linear-in-$n$ shape.
\end{itemize}
\end{hardeningbox}

\chaptersection{Tool Kit Inventory}

The machine assembled so far -- cumulative through Chapter~14. Chapters~1--13's
rows are carried forward; Chapter~14's rows are appended. Pointers only, never prose.

\begin{center}
\small
\renewcommand{\arraystretch}{1.25}
\begin{tabular}{@{}p{0.22\linewidth} l c p{0.26\linewidth} p{0.18\linewidth}@{}}
\toprule
\textbf{Tool} & \textbf{Since} & \textbf{Level} & \textbf{Trigger phrase} & \textbf{Where} \\
\midrule
\multicolumn{5}{l}{\itshape Chapters 1--13 -- carried forward} \\
Binary-tree height/size recursion & 1 & \lv1 & ``compute at the node from its two children'' & \chref{c01} \\
Bottom-up DP pattern & 2 & \lv2 & ``a node folds its children's answers'' & \chref{c02} \\
Parent-guarded rooting / representation & 4 & \lv2 & ``root a general tree with a parent parameter'' & \chref{c04} \\
Measures DFS ($\depth/\height/\size$) & 12 & \lv2 & ``how deep / how tall / how many below'' & \coderef{c09-dfs3} \\
Iterative two-phase DFS & 12 & \lv3 & ``the chain may be $5\cdot10^{5}$ deep'' & \coderef{c09-iter} \\
DP diameter (join two arms at peak) & 12 & \lv3 & ``diameter that must carry weights / a predicate'' & \thmref{c09-dpjoin} \\
Subtree tuple fold & 13.1 & \lv3 & ``an answer per node from a fixed-size tuple'' & \coderef{c10-fold} \\
In/out pair (MIS / vertex cover) & 13.2 & \lv4 & ``take or leave $v$; forced out on children'' & \coderef{c10-inout}, \thmref{c10-mis-tree} \\
Peak-harvest identity & 13.3 & \lv4 & ``sum over all pairs with an additive $F$'' & \propref{c10-pairtop} \\
Distance-$k$ histogram merge & 13.3 & \lv4 & ``count pairs at distance $k$ in one DFS'' & \coderef{c10-peak} \\
Depth-indexed state (Shape A) & 13.4 & \lv4 & ``per-node histogram up to $k$'' & \coderef{c10-distk-hist} \\
Backtracking + monotone accumulator (Shape B) & 13.4 & \lv4 & ``offline per-depth aggregate with invertible $\Phi$'' & \coderef{c10-depth-hist} \\
\midrule
\multicolumn{5}{l}{\itshape Chapter 14 -- new this chapter} \\
Capped child-merge / tree knapsack & 14.1 & \lv4 & ``state has an index $j$; merge children one at a time'' & \coderef{c27-knap}, \thmref{c27-amortize} \\
Rollback profile DFS & 14.2 & \lv4 & ``growing profile along a root-to-node path'' & \coderef{c27-profile}, \thmref{c27-profile-bound} \\
Deepest-uncovered leaf-prune greedy & 14.3 & \lv4 & ``cover a tree objective with the fewest attachments'' & \coderef{c27-greedy}, \thmref{c27-exchange} \\
Two-flag win/lose DP & 14.4 & \lv4 & ``Boolean game verdict, $k$-round parity'' & \coderef{c27-455b-soln} \\
Grundy / Hackenbush fold & 14.4 & \lv4 & ``impartial game, disjoint sum by XOR'' & \coderef{c27-grundy}, \thmref{c27-grundy-xor}, \thmref{c27-hackenbush} \\
Supply/demand 2-index tree DP (Summit) & Summit & \lv5 & ``distance-$k$ dominating sets, count'' & \coderef{c27-summit}, \probref{c27-summit} \\
\bottomrule
\end{tabular}
\end{center}

\chaptersection{Related Topics}

\begin{itemize}[nosep, leftmargin=1.4em]
  \item \chref{c01}, \chref{c02} -- backward: the capped child-merge of
        14.1 is Ch.~1's height recursion (\thmref{c01-height-recursion})
        promoted to an array state with a budget; Ch.2's bottom-up pattern
        (\chref{c02}) is the same shape with a numeric index instead of a
        scalar. Cited, never re-taught (\defref{c27-knap-state}).
  \item \chref{c04} -- backward: rooting by an explicit parent parameter
        is Ch.~4's; every Ch.14 listing inherits it verbatim.
  \item \chref{c09} -- backward: Ch.12's $\depth/\height/\size$ measures
        are the running example $T$ on which every Ch.14 hand trace sits;
        Ch.12's \thmref{c09-dpjoin} is the $\Phi=\max$ special case of
        \thmref{c27-exchange}'s coverage-superset argument at $R = \infty$.
  \item \chref{c10} -- backward: Ch.13's \thmref{c10-subtree-template} is
        the $V = 1$ special case of \thmref{c27-amortize}; Ch.13's
        \propref{c10-pairtop} peak-harvest is the merge of \secref{c27-knapsack}
        at $\Phi = \sum$ and $V = 1$; Ch.13's \secref{c10-twostate} is
        the two-flag Boolean game of \probref{c27-game-solved} in a
        different guise. This chapter pays the $\bigO{nV}$ theorem Ch.13
        previewed but did not prove (\thmref{c27-amortize}).
  \item \chref{c20} -- forward: rerooting \emph{as a theory} (any
        edge-transfer rule; answer for every root) is Chapter~15's; the
        ``rooted supply/demand'' state of \probref{c27-summit} has an
        every-root version that is Chapter~15's machinery applied to a
        growing state.
  \item \chref{c11} -- forward: per-pair \emph{query} distance and the
        $\bigO{\log n}$ LCA machinery are Chapter~16's; the mixed
        \chref{c11}-flavour entries (Duff in the Army, Lorenzo Von
        Matterhorn) use LCA on top of a Ch.14 profile DP and are the
        natural hand-off.
  \item \chref{c28} -- forward: DSU-on-tree / small-to-large upgrades
        the $\bigO{nV}$ merge to $\bigO{n \log^2 n}$ when $V$ degrades to
        $\Theta(n)$; Ch.14 previews this at \thmref{c27-amortize}'s
        complexity note but the amortised re-merge is Ch.~17's.
  \item \chref{c18} -- forward: centroid decomposition gives a
        $\bigO{n \log n}$ alternative to the $\bigO{n k^2}$ Summit DP
        when $k$ is not small and only extremal coverage is required.
\end{itemize}

\begin{lastword}[Last Word -- Chapter 14]
\emph{``Chapter~13 folded a tuple; Chapter~14 pays for a state that grows.
The root is fixed, the state has an index, and one DFS with a capped
merge is enough.''}

\medskip
The whole chapter is one discipline repeated four times: \emph{name a
growing state, name a capped merge, run one DFS}. An array per node
returns CF~855C's three-class colourings (\thmref{c27-amortize},
\probref{c27-knap-solved}); a rollback profile per start returns CF~490F's
longest increasing path (\thmref{c27-profile-bound}, \probref{c27-profile-solved});
a deepest-first sweep returns CF~1029E's leaf cover
(\thmref{c27-exchange}, \probref{c27-greedy-solved}); a two-flag
Boolean returns CF~455B's $k$-round verdict (\probref{c27-game-solved});
a Grundy integer with an XOR fold returns Hackenbush wins
(\thmref{c27-grundy-xor}, \thmref{c27-hackenbush}). The Summit is the
same shape with two indices -- supply and demand -- and a reset rule that
is Ch.13's peak-harvest identity in a new dress (\probref{c27-summit}).

\medskip
Carry three reflexes forward:
\begin{enumerate}[nosep, leftmargin=1.6em]
  \item \emph{Cap every loop by the running size.} \thmref{c27-amortize}
        is not a nicety -- forgetting $\min(\mathit{sz}, V)$ turns a
        linear-in-$n$ DP into a quadratic one, and the WA signature is
        ``TLE on star-shaped maximal tests'' (\coderef{c27-knap}).
  \item \emph{Every growing state needs a bounded index.} Profile DPs
        (\secref{c27-growing}) and supply/demand DPs (Summit) are honest
        only because $k$ and $V$ are polynomially small; when they are not,
        hand off to Ch.17's small-to-large or Ch.24's centroid
        decomposition (\thmref{c27-profile-bound}).
  \item \emph{Games and knapsacks share one skeleton.} The operator
        changes ($\max$, $\oplus$, OR-of-NOT, mex) and the reset condition
        changes, but the fold is ``one DFS, one per-frame state, children
        merged one at a time'' in all four sections. Naming the operator
        is the whole decision (\coderef{c27-kit}).
\end{enumerate}

\medskip
And three habits, drilled until boring: check that the merge's inner bounds
are $\min(\cdot, V)$ before you ship -- the theorem pays only if you cap;
compute the Grundy / win-flag values on the running tree $T$ by hand before
trusting a code answer (the four traces of this chapter all use $T$ and all
were re-derived before the listing was written); and when a problem has both
a supply and a demand, name both coordinates in the state before reaching for
a single-index DP (\probref{c27-summit}).

\medskip
\textbf{Next:} \chref{c20} takes the fixed-root DP of Chapters~13--14 and
asks the harder symmetry question -- \emph{what changes when the answer must
hold for every root at once?} Ch.~14 hands Ch.~15 a knapsack with an index,
a game with an operator, and a running tree; Ch.~15 gives you an edge-transfer
rule, a two-DFS rerooting template, and the ``parent-side piece'' this chapter
deliberately does not maintain. The analysis has been the point up to here;
from here, the root itself becomes a variable.
\end{lastword}
